% ============================================================================
% Section 2: Axioms — From Phenomenal Properties to Physical Postulates
% ============================================================================

\section{From phenomenal properties to physical postulates}
\label{sec:axioms}

Like IIT~\cite{iit4}, UHM begins with properties of experience and translates them into
mathematical postulates. Unlike IIT, UHM's postulates are grounded in a specific algebraic
structure (octonions) and a specific combinatorial structure (Fano plane), which uniquely
determine the dimensionality and architecture of the theory. UHM rests on \textbf{four axioms}
(A1--A4) and the Page--Wootters constraint of a fifth (A5); the clock register of A5
follows from A1--A4 via the spectral-triple construction (T-87, steps 1--3~\status{T}),
the constraint itself does not. These five inputs are postulates~\status{P}, each
independent; several are supported by theorems (Section~\ref{subsec:premises}). The
earlier statement that all axioms are derived as theorems is withdrawn
(Section~\ref{subsec:cohesive-closure}).

\subsection{The single primitive}

\begin{axiombox}[Axiom of Omega]
The only primitive of the theory is the $\infty$-topos of sheaves on the site of density
matrices:
\begin{equation}
  \mathfrak{T} := \bigl(\ShInf(\mathcal{C}),\; J_{\mathrm{Bures}},\; \omega_0\bigr)
\end{equation}
where $\mathcal{C} = \mathbf{DensityMat}(\mathbb{C}^7)$ is the category of $7 \times 7$
density matrices, $J_{\mathrm{Bures}}$ is the Grothendieck topology induced by the Bures
metric, and $\omega_0$ is the subobject classifier.
\end{axiombox}

The quantum formalism --- density matrices and CPTP maps, (QG) --- is posited with
the site, not derived (``why quantum theory'' stays external, T-188). Spacetime,
gravity and the structure of consciousness are derived from the primitive, the
physical readings of spacetime and matter through named bridge premises
(Section~\ref{subsec:premises}). The coherence matrix $\Gam$ is an object of this category:
$\Gam \in \mathrm{Ob}(\mathcal{C})$.

\subsubsection*{The four axioms}

\begin{enumerate}[leftmargin=2em]
  \item \textbf{(A1) Structure \status{P}:} Reality is an $\infty$-topos $\ShInf(\mathcal{C})$ over
        the site of density matrices $\Dens(\mathbb{C}^7)$, generating continuous CPTP
        dynamics $\Lind_\Omega$ (Lindbladian form via GKSL theorem).
  \item \textbf{(A2) Metric \status{P}:} The Grothendieck topology $J_{\mathrm{Bures}}$ is induced
        by the Bures distance:
        \begin{equation}
          d_B(\Gam, \Gam') = \sqrt{2 - 2\sqrt{F(\Gam, \Gam')}}
          \label{eq:bures}
        \end{equation}
        where $F(\Gam, \Gam') = \bigl(\Tr\sqrt{\sqrt{\Gam}\,\Gam'\sqrt{\Gam}}\,\bigr)^2$
        is the Uhlmann fidelity. By the Petz classification~\cite{petz1996}, the Bures metric is the
        \emph{minimal} monotone Riemannian metric on $\Dens(\Hil)$
        (in the classical case, the Fisher metric is unique by
        Chentsov's theorem~\cite{chentsov1982}; in the quantum case,
        infinitely many monotone metrics exist, but Bures is minimal).
        This choice ensures $\Pcrit = 2/7$ is the least restrictive threshold.
        \emph{Support:} T-187~\status{T} proves that Bures is the uniquely canonical
        Petz-enrichment, fixed by three independent characterizations (Petz extremality,
        Uhlmann purification, SLD-Fisher saturation), and T-189~\status{T} recasts it as a
        maximum-entropy selector (Section~\ref{subsec:cohesive-closure}); the axiom
        itself stays a postulate.
  \item \textbf{(A3) Dimension \status{P}:} The base Hilbert space dimension is $N = 7$, grounded
        in octonion structure $G_2 = \mathrm{Aut}(\Oct)$ and the Fano plane $\Fano$
        (Section~\ref{subsec:why-seven}). $N \geq 7$ is a theorem (Theorem~S~\status{T});
        strict necessity of $N = 7$ is \status{C at (P1$_6$)}.
  \item \textbf{(A4) Scale \status{P}:} Characteristic frequency
        $\omega_0 > 0$ relates internal time $\tau$ to physical time $t$. A system with
        $\omega_0 = 0$ has no dynamics, hence no regeneration, hence is not viable; the
        \emph{value} of $\omega_0$ is a free parameter and differs between holons.
\end{enumerate}

The fifth axiom is stratified (T-87):

\begin{enumerate}[leftmargin=2em, start=5]
  \item \textbf{(A5) Page--Wootters clock and constraint~\cite{page1983}:} the clock
        register --- the tensor factor $\mathcal{H}_O = \mathbb{C}[\mathbb{Z}_7]$ --- follows from
        A1--A4 (T-87, steps 1--3~\status{T}); the constraint $\hat{C}\,\Gam_{\mathrm{total}} = 0$,
        the support condition $\mathrm{supp}\,\Gam \subseteq \ker\hat C$, is an independent
        postulate~\status{P} (stationarity of a mixed state gives only
        $[\hat{C}, \Gam_{\mathrm{total}}] = 0$; the earlier ``A5 is a consequence of A1--A4'' is
        withdrawn). Emergent time from the O-dimension is developed in
        Section~\ref{subsec:emergent-time}.
\end{enumerate}

\subsubsection*{Premises in use}
\label{subsec:premises}

The independent axiomatic content is A1--A4 plus the constraint of A5: five inputs,
each independent (a state off $\ker\hat C$ satisfies A1--A4 and violates the constraint).
On top of them the current corpus uses, as listed on the premises page of the archive:
\begin{itemize}[leftmargin=2em,nosep]
  \item \textbf{Bridge premises of physics \status{H}:} (Cl$_0$) --- fermions are vectors
        of the spinor module $\mathcal S = \mathbb C\otimes\Oct$ of the octonionic Clifford
        system; (P) --- the tangent vectors of spacetime are the Hermitian forms on the
        spinor factor of the fermion field, with a causal form preserved by every
        internal-structure-preserving transformation ((P)~$\Leftrightarrow$~(L)~$\wedge$~(W),
        Theorem~48e). Neither is derivable from the axioms, and the two are independent.
        T-347~\status{T} (as mathematics; no status changes) closes the routes from the
        holon: (Cl$_0$)~$\Leftrightarrow$~(Mod), the holon's product acting on matter --- a
        real-linear $\rho:\Oct\to\mathrm{End}_{\mathbb R}(F)$ with $\rho(1) = 1$ and
        $\rho(x)^2 = \rho(x^2)$, commuting with $i$; no property of the holon (viability,
        the window, the self-model, regeneration, $G_2$-rigidity) decides (Cl$_0$), (P) or
        (W$_0$), because everything built from $\Gam$ is tensorial while matter is spinorial
        (the $2\pi$-rotation is $-1$ on $\Oct$ and $+1$ on $\mathbb{R}^7$); the maximality of
        $\mathrm{Spin}(9)$ presupposes (Mod), and the minimal spinor factor ($n = 1$) has no
        spatial direction.
  \item \textbf{Principle of the self-model \status{Pr}:} (Max$\Phi$) --- the anchor of the
        self-model is maximally integrated, $\Phi(\rho_a) = 6$ (T-334).
  \item \textbf{Strict necessity of $N = 7$ \status{H}:} (P1$_6$); $N \geq 7$ needs no premise.
  \item \textbf{Free parameters:} the associator coupling $\kappa$ of $V_{\mathrm{Gap}}$ (free:
        no derived source fixes it, T-331), $\mu^2$, $\lambda_4$, the regeneration rate (fixed by none of five routes:
        threshold, extremum, criticality, normalisation, coarse-graining; T-346), the
        Fano weight $\alpha$, $\omega_0$, and the strong-CP angle $\bar\theta$ (the strong CP
        problem is open \status{Pr}; confinement Theorem~3.1c).
  \item \textbf{Identification hypotheses \status{H}:} the sector values (SV) of the Gap
        vacuum, the generation reading (GC) in broken form, (UP) at leading order,
        (FE) and (SA) in the axis frame, the Higgs identification, T-186(a), and the
        reconstruction and aperiodic-clock conditions of T-119/T-120 read as physics;
        (HOL) is an interpretation~\status{I}.
\end{itemize}
Results labelled ``\status{T} as mathematics'' use only the axioms; their physical
readings carry the premises above. The former inputs (Alt), (MP), (MM), (Q), (RT),
(Col), (Pure), (AGG) and (ND) are discharged.

\paragraph{Principle of Informational Distinguishability (PID, T-16~\status{O}).}
Given A1 and A2, a further principle follows tautologically: two states are ontologically
distinct if and only if they are distinguishable by $J_{\mathrm{Bures}}$-coverings. This
connects the mathematical formalism to consciousness: subjective distinctness and physical
distinguishability are the same relation. PID is not an additional axiom --- it is a
definition ~\status{O} (T-16 provides the tautological consequence under A1+A2).

The Bures metric on $\Dens(\mathbb{C}^7)$ defines a topology in the point-set sense. This topology generates a Grothendieck site: coverings are families of open inclusions that jointly cover an open set. By Lurie's theorem (HTT 6.2.2.7), the $\infty$-category of sheaves on this site is an $\infty$-topos, provided the Giraud axioms (descent, universal colimits, disjoint coproducts, effective groupoid objects) are satisfied. For the site of density matrices with the Bures topology, these axioms hold by construction~\cite{lurie2009}. The site axioms (identity, stability under pullback, transitivity) are verified via CPTP contractivity of the Bures metric (Uhlmann 1976) and the essentially small presentation of the compact metrizable $\Dens(\mathbb{C}^7)$ (Johnstone, \emph{Elephant} C2.1.9--12).

\textbf{Why Bures specifically?} T-187~\status{T} provides \emph{three logically
independent} mathematical witnesses plus \emph{one physical recasting}: (Char-I)~Petz
extremality (minimal in the Petz poset), (Char-II)~Uhlmann purification universality,
(Char-III)~SLD-Cram\'{e}r--Rao saturation, plus (Char-IV, T-189~\status{T})~a
\emph{physical selector} $g_{ij} = \tfrac14 C^{\mathrm{SLD}}_{ij}$ matching the metric
against the Petz-free SLD bilinear form (defined from
$\partial\rho = \tfrac12(L_i\rho + \rho L_i)$ without reference to any metric). Char-IV
reduces to Char-III via $\mathcal F_{\mathrm{SLD}} = 4g_B$ (Braunstein--Caves 1994; the
earlier identity $g_B^{-1} = C_{\mathrm{SLD}}$ is false and corrected), so it is
not a fourth logically independent witness --- it is a statistical-mechanical recasting
that strengthens the physical motivation of T-187 without adding mathematical redundancy.
T-187 retains status \status{T} on the strength of Char-I (Petz extremality) alone.

\textbf{Why $\infty$-topos?}~\cite{lurie2009} An $\infty$-topos provides three things simultaneously:
(1)~an internal logic (Heyting algebra, not Boolean --- capturing the fact that propositions
about consciousness need not satisfy the law of excluded middle);
(2)~a notion of ``space'' (Bures metric on density matrices, from which physical spacetime
emerges); (3)~higher morphisms at all levels (allowing self-reference without paradox,
via Lawvere's fixed-point theorem). No simpler structure provides all three.

\subsection{Why seven dimensions: the Fano plane and octonions}
\label{subsec:why-seven}

The number 7 is not chosen: $N \geq 7$ is \emph{derived} (Theorem~S~\status{T}), and
strict necessity of $N = 7$ holds \status{C at (P1$_6$)}, P1 for a competing
decomposition. Two complementary
arguments fix it: a \emph{functional} count that determines the number, and an
exceptional \emph{algebraic} structure that determines the shape of the
seven-dimensional system (the octonions and the Fano plane are two faces of
one object: the Fano plane is the incidence geometry of the octonion
multiplication table).

\paragraph{Argument 1: Octonion algebra (the structure).}
The octonions $\Oct$ form the unique normed division algebra in 8 dimensions (by the
Hurwitz theorem~\cite{hurwitz1898}, the only normed division algebras are $\mathbb{R}$,
$\mathbb{C}$, $\mathbb{H}$, and $\Oct$). The imaginary part of $\Oct$ has dimension 7:
$\dim(\mathrm{Im}(\Oct)) = 7$. The automorphism group of $\Oct$ is the exceptional
Lie group $G_2$, which has dimension 14.

\paragraph{Argument 2: Fano plane $\Fano$ (the same structure, incidence face).}
The Fano plane is the smallest finite projective plane: 7 points and 7 lines,
where each line passes through exactly 3 points, and each point lies on
exactly 3 lines (Figure~\ref{fig:fano}). It is the incidence geometry of
the octonion multiplication table.

\begin{figure}[H]
\centering
\input{figures/fano-plane}
\caption{\textbf{The Fano plane $\Fano$ --- the architecture of self-observation.}
Seven points represent the seven dimensions of $\Gam$: A~(Articulation), S~(Structure),
D~(Dynamics), L~(Logic), E~(Interiority), O~(Ground), U~(Unity). Seven lines represent
the seven triples through which the self-modeling operator $\vp$ observes the system.
Each dimension participates in exactly 3 triples; each triple binds exactly 3 dimensions.
This ``maximally democratic'' structure is unique at 7 points and does not exist at 5 or 6.}
\label{fig:fano}
\end{figure}

\paragraph{Why the Fano plane matters for consciousness.}
Self-observation --- the act of a system modeling itself --- requires examining correlations
between dimensions. A naive approach (examining all pairs) gives $\binom{7}{2} = 21$
channels, which is redundant and destroys coherence. The Fano plane provides the
\emph{minimal} set of channels (7 triples) that covers every pair exactly once. This is
the mathematical content of ``balanced self-observation'': the system observes itself
through every possible triple of dimensions simultaneously, preserving enough coherence
to survive.

\paragraph{Argument 3: Diagnosability (the selector --- and its identity with the algebra).}
A third argument fixes $N = 7$ from an entirely different direction --- fault diagnosis ---
and it turns out to be the \emph{same} fact as Argument~1. Demand that the seven-channel
frame be \emph{perfectly single-fault diagnosable}: every single-channel fault is uniquely
located by a fixed family of parity checks, with no syndrome wasted. Perfect binary
single-error-correcting codes exist only at lengths $2^r - 1$; with two mild nondegeneracy
and rigidity conditions this forces $r = 3$, hence $N = 7$ with the Hamming/Fano incidence,
uniquely (Theorem~$\Sigma$, T-224~\status{T}). This is not a third coincidence with the
octonionic Argument~1 but, in part, a proven identity. A frame is perfectly diagnosable
\emph{iff} it is anisotropic --- carries no ``null'' channel on which a fault would be
silent (an idempotent zero-divisor) --- and anisotropy of a normed algebra is exactly the
division property (T-244~\status{T}). Hence ``the last normed division algebra'' (Hurwitz)
and ``the last perfectly self-diagnosing frame'' (Theorem~$\Sigma$) name the same seven
channels; the pivotal sign $e_g^2 = -1$ on each channel --- the quaternionic
Frobenius--Schur indicator, gauge-invariant --- is forced by the same nondegeneracy
(T-243~\status{T}). Three considerations thus fix the value: algebraic structure
(Arguments~1--2, the octonion algebra and its Fano incidence face), diagnostic selection
(Argument~3), and the functional minimality proved immediately below (T-113~\status{T}).
They do not merely converge on seven --- algebraic structure and diagnostic selection are
literally one theorem (anisotropy $=$ division), so the value is overdetermined rather than
fitted.

\begin{theorem}[Minimality of 7 for autonomous cognition, T-113 \status{T}]
For $N < 7$, no finite projective plane with the required incidence properties exists
(the Fano plane $\Fano$ is the unique projective plane of order 2, and smaller orders
are impossible). Consequently, a system with fewer than seven functionally independent
dimensions cannot simultaneously support self-observation, regeneration, and autonomous
learning (T-113, T-114~\status{T}): the optimal number of learning steps diverges,
$n^\ast = \infty$. $N = 7$ is the mathematical minimum for autonomous consciousness.
\end{theorem}

\subsection{The seven dimensions}
\label{subsec:seven-dimensions}

The coherence matrix $\Gam$ is a $7 \times 7$ Hermitian positive semidefinite matrix with
$\Tr(\Gam) = 1$. Its seven diagonal elements $\gamma_{kk}$ represent the ``weight'' of
each dimension; its off-diagonal elements $\gamma_{ij}$ ($i \neq j$) represent coherences
(quantum correlations) between dimensions.

\begin{table}[H]
\centering
\caption{The seven dimensions of the coherence matrix $\Gam$.}
\label{tab:dimensions}
\begin{tabular}{@{}clll@{}}
\toprule
\textbf{Index} & \textbf{Dimension} & \textbf{What it governs} & \textbf{Fano lines} \\
\midrule
$k=1$ & \textbf{A} (Articulation) & Perception, sensory bandwidth & $\{A,S,L\}$, $\{E,U,A\}$, $\{O,A,D\}$ \\
$k=2$ & \textbf{S} (Structure)    & Spatial organization, form    & $\{A,S,L\}$, $\{S,D,E\}$, $\{U,O,S\}$ \\
$k=3$ & \textbf{D} (Dynamics)     & Temporal unfolding, change    & $\{S,D,E\}$, $\{D,L,U\}$, $\{O,A,D\}$ \\
$k=4$ & \textbf{L} (Logic)        & Reasoning, abstraction        & $\{A,S,L\}$, $\{D,L,U\}$, $\{L,E,O\}$ \\
$k=5$ & \textbf{E} (Interiority)  & Subjective experience itself  & $\{S,D,E\}$, $\{L,E,O\}$, $\{E,U,A\}$ \\
$k=6$ & \textbf{O} (Ground)       & Energy, metabolism, resources  & $\{L,E,O\}$, $\{U,O,S\}$, $\{O,A,D\}$ \\
$k=7$ & \textbf{U} (Unity)        & Integration, the sense of self & $\{D,L,U\}$, $\{E,U,A\}$, $\{U,O,S\}$ \\
\bottomrule
\end{tabular}
\end{table}

\subsection{Two-aspect monism: reframing the hard problem}
\label{subsec:two-aspect}

\begin{figure}[H]
\centering
\input{figures/two-aspect}
\caption{\textbf{Two-aspect monism.} The coherence matrix $\Gam$ is a single object with
two inseparable projections. The external projection ($\mathrm{Map}_{\mathrm{ext}}$) yields
physical observables --- dynamics, structure, spacetime geometry. The internal projection
($\mathrm{Map}_{\mathrm{int}}$, restricted to the E-sector) yields phenomenal content ---
qualia, the qualitative character of experience. There is no bridge between them because
they are not separate.}
\label{fig:two-aspect}
\end{figure}

The hard problem presupposes that physics and experience are ontologically
distinct and asks how they interact. UHM reframes this presupposition
(Figure~\ref{fig:two-aspect}):

\begin{definition}[Two-aspect monism \status{T}/\status{O}]
\label{def:two-aspect}
$\Gam$ is a single ontological primitive. Its external aspect
($\mathrm{Map}_{\mathrm{ext}}(\Gam)$) is physics; its internal aspect
($\mathrm{Map}_{\mathrm{int}}(\Gam)$) is experience. The phenomenal functor
$F: \mathbf{Phys} \to \mathbf{Phen}$ is an \emph{isomorphism} between the physical
and interiority categories --- preserving the structure of morphisms (as proclaimed by
Spinoza's E2P7). The functor $F$ is unique given the axioms (theorem of uniqueness of
the phenomenal functor~\status{T}); its identification with the infinitesimal flat
modality of cohesion (T-186(a), \S\ref{subsec:cohesive-closure}) is a
hypothesis~\status{H}. The functorial part is \status{T}; the
ontological identification "internal aspect = experience" is \status{O}
(definitional; PH clause of the characterising properties). However, the self-modeling operator $\vp$ is
\emph{not} an isomorphism: $\mathrm{Gap}(\Gam, \vp(\Gam)) > 0$ (T-55~\status{T}, Lawvere
incompleteness~\cite{lawvere1969}). The categories are structurally identical; the
system's self-knowledge of them is structurally incomplete.

\medskip\noindent\textit{Epistemic status.} Two-aspect monism is the semantic bridge
from the mathematical primitive $\Gam$ to the phenomenology of experience. The functorial
structure (that $F$ exists and is unique given the axioms) is a theorem~\status{T}. The
ontological claim that the internal aspect \emph{is} experience (rather than merely
\emph{models} it) is the \emph{defining property} (PH) of a viable holon, read from
axiom $\Omega^7$ as a whole --- not a sixth axiom. The Axiomatic Closure T-190 derives
A1--A5 from (AP)+(PH)+(QG)+(V)+MaxEnt only \status{C} under the constraint of A5, the
hypothesis T-186(a) and the CPTP-monotonicity of the enrichment of A2; the earlier claim of \textbf{zero} independent axioms is withdrawn
(Section~\ref{subsec:premises}).
\end{definition}

This position has a deep philosophical genealogy:
\begin{itemize}[leftmargin=2em]
  \item \textbf{Spinoza (1677):} ``The order and connection of ideas is the same as the
        order and connection of things'' (Ethics II, Prop.~7) --- the exact precursor of
        the phenomenal functor $F$.
  \item \textbf{Russell (1927):} ``There is a `neutral stuff' neither mental nor physical''
        --- $\Gam$ is Russell's neutral stuff.
  \item \textbf{Chalmers (1996):} ``There are psychophysical laws connecting physics to
        consciousness'' --- UHM: no laws needed; one object, two projections.
\end{itemize}

\paragraph{The categorical gap.}
While the phenomenal functor $F$ is an isomorphism between categories (physical structures
map bijectively to phenomenal structures), the self-modeling operator $\vp$ is \emph{not}:
$\mathrm{Gap}(\Gam, \vp(\Gam)) > 0$ (T-55~\status{T}). This means the system can never
fully model \emph{itself} --- the external description and internal experience are
structurally identical, but no finite system can access this identity from within.
The explanatory gap is not between physics and experience (they are one); it is between
the system and its own self-knowledge. This is proved via Lawvere's fixed-point
theorem --- the same phenomenon as G\"{o}del incompleteness, applied to self-observation.

\paragraph{Gap as source of meaning.}
The Gap is not merely an epistemological limitation --- it is the \emph{source} of
semantic content. The Gap operator $\hat{\mathcal{G}} := \mathrm{Im}(\Gam)$ is an
antisymmetric matrix in the Lie algebra $\mathfrak{so}(7)$ with spectrum
$\{\pm i\lambda_1, \pm i\lambda_2, \pm i\lambda_3, 0\}$~\status{T}. Its Frobenius
norm $\mathcal{G}_{\mathrm{total}} = 2(\lambda_1^2 + \lambda_2^2 + \lambda_3^2)$
measures total opacity between external and internal descriptions.

Within the spectral triple formalism, Gap exactly equals the curvature of the connection
on the internal noncommutative geometry~\status{T}: zero Gap means flat connection
(external fully determines internal); nonzero Gap means curvature (under cyclic changes
of external parameters, the internal state acquires a geometric phase --- a Berry phase
analogue). The richer the phenomenal content, the greater the curvature.

The Gap decomposes as $\hat{\mathcal{G}} = \hat{\mathcal{G}}_{G_2} +
\hat{\mathcal{G}}_\perp$ into a $G_2$-preserving part (14-dimensional, structurally
healthy) and a $G_2$-breaking part (7-dimensional, potentially pathological)~\status{T}.
At full opacity rank ($r = 3$, generic spectrum), the configuration is
\emph{topologically protected}: the flag variety $G_2/T^2$ has
$\pi_2(G_2/T^2) \cong \mathbb{Z}^2$ (the rank-$2$ weight lattice), so nontrivial
second-homotopy classes obstruct continuous deformation of a rank-$3$ Gap to the
trivial (zero-Gap) state (T-69: $\pi_2(G_2/T^2) \cong \mathbb{Z}^2$~\status{T}; the
energy barriers $\geq 6\mu^2$ protecting the vacuum are Hessian eigenvalues of the
retracted sector parametrisation of T-64, so the protection is \status{C at (SV)}).
A sufficiently rich inner life is in this sense topologically stable.

\subsection{Cohesive closure: what survives (T-185--T-190)}
\label{subsec:cohesive-closure}

\paragraph{Attribution.} The cohesive-closure layer of UHM (theorems T-185 and beyond)
rests on a specific mathematical lineage: Lawvere's axiomatic cohesion~\cite{lawvere1973}
(1970s), refined and extended to the differential-cohesive $\infty$-topos framework by
Schreiber~\cite{schreiber2013} (DCCT, 2013--2016), with further super-cohesive
developments~\cite{sati-schreiber2018}. The cohesion-dependent statements that follow
--- T-185 (modalities), T-186 (phenomenal functor $F \cong \&|_{\mathcal D}$), T-188
(hard-problem localization) --- use Schreiber's modal calculus as their structural
scaffolding. The pre-T-182 core of UHM (T-1 through T-181) stands independently on
octonionic, Fano, Lindbladian, and Bures foundations. Two former members of this layer
no longer use cohesion: the U-projection is the $G_2$-twirl
$X \mapsto \tfrac17\Tr(X)\,\mathbf 1$ (T-212$'$~\status{T}), and its identification with
the rheonomy modality $\mathrm{Rh}$ is retracted~\status{$\times$} ($\mathrm{Rh}$ preserves
global points and is not among the modalities of differential cohesion); T-221 is the
relationalist route through the List/DeBrota no-go results (\S\ref{subsec:foundational-closure}).

Theorem T-185 is stratified. (i)~In any differentially cohesive $\infty$-topos the
modalities $\Pi \dashv \flat \dashv \sharp$ and $\mathrm{Red} \dashv \Pi_{\inf} \dashv
\flat_{\inf}$ exist~\status{T}. (ii)~That the site $(\mathbf{DensityMat}, J_{\mathrm{Bures}})$
is differentially cohesive is an assumption~\status{C}: it has no $\infty$-cohesive site,
and the petit $\infty$-topos of the space $\Dens(\mathbb{C}^7)$ is not cohesive at all.
(ii$'$)~$\Dens(\mathbb{C}^7)$ with its rank strata is an object of the differentially
cohesive $\mathrm{SynthDiff}\infty\mathrm{Grpd}$, the strata $\Dens_k$ are submanifolds of
dimension $14k - k^2 - 1$ with shapes $\Pi\Dens_k \simeq \mathrm{Gr}_k(\mathbb{C}^7)$, and
the seven modalities $\mathrm{Id}, \Pi, \flat, \sharp, \mathrm{Red}, \Pi_{\inf}, \flat_{\inf}$
are pairwise distinct~\status{T}. Their assignment to the seven dimensions is an
interpretation~\status{I}; the former list with $\mathrm{Rh}$ in place of $\mathrm{Red}$ is
retracted~\status{$\times$}.

For a \emph{stable} coefficient object $E$ of a cohesive $\infty$-topos
$\mathbf{H}$~\cite{schreiber2013} there is a canonical \emph{differential cohomology
hexagon} (DCCT v1, Prop.~4.1.17), whose structural maps include
\begin{equation}
  \flat(E) \to E \to \flat_{\mathrm{dR}}(E).
  \label{eq:hexagon}
\end{equation}
For stable $E$ it forces the relationship between the internal aspect ($\flat$, flat
modality), the external structure ($\Pi$, shape modality), and the curvature
($\flat_{\mathrm{dR}}$, de Rham coefficient). An earlier version of this section stated the
hexagon for \emph{any} coefficient object; that is retracted --- it does not apply to
$\mathbf{B}G_2$, which is not stable.

\begin{theorembox}[T-186: Cohesive Closure --- (a) \status{H}; (b), (c) \status{$\times$}]
Let $\mathfrak{T} = \ShInf(\Dens(\mathbb{C}^7), J_{\mathrm{Bures}})$ be the UHM
$\infty$-topos with differentially cohesive structure (assumed in the original statement;
available in the corrected form T-185 (ii$'$)). As originally stated:

\textbf{(a) Phenomenal necessity.} The phenomenal functor $F$ is naturally isomorphic
to the infinitesimal flat modality:
$F \cong \&\big|_{\Dens}$,
and the Postnikov filtration of $\&(\Gam)$ reproduces the interiority hierarchy L0--L4.
\emph{Status: hypothesis}~\status{H}. The modality $\&$ exists on the state space
(T-185 (ii$'$)), but the natural isomorphism is asserted, not constructed; modal
constructions commute with every unitary conjugation of $\Dens(\mathbb{C}^7)$ and cannot
pick out the E-axis of the pointer frame, which $F$ uses.

\textbf{(b) Exact time.} ``The Page--Wootters conditional states are sections of
$\flat(\Gam_{\mathrm{total}})$, governed exactly by the counit, with no
$O(H_{\mathrm{int}})$ correction.'' \emph{Retracted}~\status{$\times$}: with a clock--system
interaction the conditional dynamics is time-nonlocal (Smith--Ahmadi 2019), relative to a
seven-periodic clock every conditional dynamics is periodic, and the counit carries no
information about $H_{\mathrm{int}}$.

\textbf{(c) Unconditional $\Lambda > 0$.} ``$\Delta F = \|\mathrm{curv}(\Gam)\|^2 =
\omega_0^2\,\mathcal{G}_{\mathrm{total}} > 0$ via Chern--Weil and T-55.''
\emph{Retracted}~\status{$\times$}: the hexagon holds for stable coefficients only, not for
$\mathbf{B}G_2$, and every $G_2$-bundle over the contractible $\Dens(\mathbb{C}^7)$ is
trivial, so all its characteristic classes vanish. $\Delta F > 0$ and $\Lambda > 0$ return
to their conditional status~\status{C} (on the spectral details of $D_{\mathrm{int}}$).
\end{theorembox}

\paragraph{Worked example (curvature, not an unconditional $\Delta F$).}
Consider a holon at $P = 0.35$ with $\omega_0 = 40$\,Hz and three
representative coherences on Fano lines:
$\gamma_{EO} = 0.08\,e^{i\pi/3}$,
$\gamma_{AE} = 0.06\,e^{i\pi/4}$,
$\gamma_{OU} = 0.05\,e^{i\pi/5}$,
with $\mathrm{Gap}(i,j) = |\sin\theta_{ij}|$.
The total Gap is
$\mathcal{G}_{\mathrm{total}} = \sum_{i<j}|\gamma_{ij}|^2\,\mathrm{Gap}(i,j)^2
\geq 0.0048 + 0.0018 + 0.0009 = 0.0075$,
and the curvature of the Serre bundle (T-73~\status{T}:
$\|\mathrm{Curv}\|_{ij}^2 = \omega_0^2|\gamma_{ij}|^2\,\mathrm{Gap}(i,j)^2$) has squared norm
$\omega_0^2\,\mathcal{G}_{\mathrm{total}} \geq 1600 \cdot 0.0075 = 12.0$.
Reading this number as a free-energy gradient $\Delta F > 0$ was part~(c) of T-186 and is
retracted; the free-energy fuel for regeneration is conditional~\status{C}.

\begin{theorembox}[T-187: Canonicity of the Bures Enrichment \status{T}]
Within the Petz family of CPTP-monotone Riemannian metrics on $\Dens(\mathbb{C}^7)$, the
Bures metric $d_B$ is uniquely fixed by three independent characterizations:

\textbf{(Char-I) Petz extremality:} $d_B$ is the pointwise minimum of the Petz poset
($d_B \le d_f$ for all Petz $f$), equivalently the terminal object of the Petz diagram
in Lawvere-enriched $[0,\infty]$-categories~\cite{lawvere1973}.

\textbf{(Char-II) Uhlmann universality:} $d_B$ is the unique metric satisfying
$d_B(\rho,\sigma) = \inf_{|\psi\rangle,|\varphi\rangle}\||\psi\rangle - |\varphi\rangle\|$
over purifications (Uhlmann 1976).

\textbf{(Char-III) SLD-Fisher:} $4g_B$ = the symmetric-logarithmic-derivative Fisher metric,
uniquely saturating the quantum Cram\'{e}r--Rao bound (Braunstein--Caves 1994~\cite{braunstein1994}).

All three select the same $d_B$; any other Petz metric gives the same classical
$\infty$-topos (compact bi-Lipschitz equivalence at the topological level) but a different
enrichment. The Grothendieck topology $J_B$ is generated by the $\varepsilon$-$\delta$
coverage (Johnstone, \emph{Elephant} C2.1.10); transitivity is automatic from the generation.
This supports A2; the axiom itself is listed among the postulates
(Section~\ref{subsec:premises}).
\end{theorembox}

\paragraph{Axiomatic closure is conditional (T-190).}
With T-187, T-189 and T-87, the axioms A1--A5 are derivable from the characterising
properties $(\text{AP})+(\text{PH})+(\text{QG})+(\text{V})+\text{MaxEnt}$ of a viable holon
only under three conditions: the Page--Wootters constraint of A5 is assumed, the route
to A1 through T-186(a) is a hypothesis, and the enrichment of A2 is CPTP-monotone,
$d(\mathcal E\rho,\mathcal E\sigma) \le d(\rho,\sigma)$ for every CPTP map $\mathcal E$. The
third condition was named on 2026-09-26 (Lemma~M of the corpus, \status{T}): it follows from
none of (AP), (PH), (QG), (V), MaxEnt --- the Hilbert--Schmidt metric, in which (V) is
written, satisfies them all and grows by the factor $\sqrt2$ under a partial-trace channel on
$\mathbb{C}^7$ --- but it follows from the operational reading (O) of the enrichment as the
best distinguishability over measurements, which also gives Bures directly
(Fuchs--Caves~\cite{fuchs1995}). The derivations are:
\begin{itemize}[leftmargin=2em,nosep]
  \item \textbf{A1 \status{H}:} $\infty$-topos $\ShInf(\mathcal{C})$ from T-76 (Grothendieck site
        level \status{T} via Bures topology, Uhlmann 1976 CPTP contractivity;
        Exp-extension conditional on Giraud axiom verification per Claim 10.2)
        + T-186(a), a hypothesis.
  \item \textbf{A2 \status{C} (monotonicity):} the topology is forced (every continuous
        distance on the compact $\Dens(\mathbb{C}^7)$ induces the standard one); \emph{within
        the CPTP-monotone metrics} Bures is the uniquely canonical Petz-enrichment via T-187~\status{T} (triple
        characterisation: Char-I Petz extremality, Char-II Uhlmann universality, Char-III
        SLD-Fisher Cram\'er--Rao saturation) + T-189~\status{T} (Char-IV, physical recasting via
        MaxEnt: $g_{ij} = \tfrac14 C^{\mathrm{SLD}}_{ij}$, equivalently $g_B^{-1} = 4\,C_{\mathrm{est}}$);
        the monotonicity itself is the third condition (this line read ``via T-187 + T-189''
        until 2026-09-26).
  \item \textbf{A3:} $N \geq 7$ from Theorem~S~\status{T}; the \emph{structure} from the
        Bridge T-15 chain
        $(\text{AP})+(\text{PH})+(\text{QG})+(\text{V}) \xrightarrow{\text{T-15}} \mathrm{BIBD}(7,3,1) \to \mathrm{PG}(2,2) \to \mathbb O \to G_2$
        with the canonical orientation of the Fano lines (T15-canon~\status{T}; the chain
        consumes $N = 7$ from Theorem~S); strict necessity of $N = 7$ is
        \status{C at (P1$_6$)}.
  \item \textbf{A4:} $\omega_0 > 0$ forced by viability (V): $\omega_0 = 0 \Rightarrow$ no
        dynamics $\Rightarrow$ (AP) violated; its value is a free parameter.
  \item \textbf{A5:} the clock register from T-87, steps 1--3~\status{T}; the constraint
        $\hat C\,\Gam = 0$ is not derived (T-87, step 4, is \status{C} under the support
        condition of Property~2).
\end{itemize}

\begin{theorembox}[T-190: Axiomatic Closure of UHM \status{C}]
\label{thm:t190-axiomatic-closure}
A1--A5 of $\Omega^7$ are derivable from
$(\text{AP})+(\text{PH})+(\text{QG})+(\text{V})+\text{MaxEnt}$ \emph{under three conditions}:
the Page--Wootters constraint is assumed (T-87, step 4), the route to A1 through
T-186(a) is a hypothesis, and the enrichment of A2 is CPTP-monotone (Lemma~M). The
derivation chain per axiom:
A1 via T-76+T-186(a); A2 via monotonicity + T-187 + T-189 (triple Petz/Uhlmann/SLD-Fisher + MaxEnt
physical recasting); A3 via Theorem~S and Bridge T-15; A4 via (V) necessity;
the clock register of A5 via T-87, steps 1--3. The earlier statement --- all five axioms
are theorems and UHM has \textbf{zero} independent axioms --- is
retracted~\status{$\times$}: the constraint remains an independent assumption. Proof details
appear in Section~\ref{subsec:foundational-closure} (see also the companion localisation
theorem T-188~\status{C}, the $\vp$-tower convergence T-191 under backbone dominance, and the
target-category verification T-192 for $\mathbf{Exp}^{(2)}$).
\end{theorembox}
